\section{Methodology, Architecture, and Algorithms}

% ============================================================
% S8: METHODOLOGY - DSRM
% ============================================================
\begin{frame}{Methodology: Design Science Research}
    MEDF follows the Design Science Research Methodology (DSRM) by Peffers et al.~\cite{peffers2007dsrm}, providing academic legitimacy beyond pure software engineering.

    \vspace{8pt}
    \begin{enumerate}
        \item \textbf{Problem Identification:} Principles-to-practice gap in AI ethics (Chapters 2--3).
        \item \textbf{Objective Definition:} Quantifiable, multi-stakeholder evaluation (Chapter 3).
        \item \textbf{Design and Development:} Architecture, algorithms, and implementation (Chapters 6--9).
        \item \textbf{Demonstration and Evaluation:} Three real-world case studies (Chapters 10--11).
    \end{enumerate}
\end{frame}
% NOTES: DSRM positions the project as a research artifact, not just code. The six
% DSRM activities map directly onto report chapters. Approximately 60 seconds.
% Transition: "The first design decision was harmonizing three governance frameworks."

% ============================================================
% S9: FRAMEWORK HARMONIZATION
% ============================================================
\begin{frame}{Framework Harmonization: Six Unified Dimensions}
    Three governance frameworks harmonized into one evaluation model.

    \vspace{4pt}
    \begin{columns}[T]
        \begin{column}{0.35\textwidth}
            {\small
            \textbf{Source Frameworks:}
            \begin{itemize}
                \item EU ALTAI~\cite{aihleg2019ethics, aihleg2020altai}.
                \item NIST AI RMF~\cite{nist2023rmf}.
                \item Singapore MGAF~\cite{pdpc2020mgaf}.
            \end{itemize}
            }
        \end{column}
        \begin{column}{0.60\textwidth}
            {\small
            \textbf{Six Unified Dimensions:}
            \begin{enumerate}
                \item Transparency and Explainability.
                \item Fairness and Non-discrimination.
                \item Safety and Robustness.
                \item Privacy and Data Governance.
                \item Human Agency and Oversight.
                \item Accountability.
            \end{enumerate}
            }
        \end{column}
    \end{columns}

    \vspace{6pt}
    {\small Derived from overlapping themes across all three frameworks, informed by Jobin et al.~\cite{jobin2019global}, Floridi et al.~\cite{floridi2018ai4people}, and Zhang et al.~\cite{zhang2023fairness}. Stored as version-controlled YAML files.}
\end{frame}
% NOTES: Explain that the harmonization is a central contribution. Each framework
% assigns different intrinsic weights to each dimension. Approximately 60 seconds.
% Transition: "The system architecture implements this model as a three-tier platform."

% ============================================================
% S10: SYSTEM ARCHITECTURE
% ============================================================
\begin{frame}{System Architecture: Three-Tier Design}
    \begin{columns}[T]
        \begin{column}{0.40\textwidth}
            \begin{itemize}
                \item \textbf{Presentation:} Streamlit dashboard with interactive sliders and radar charts.
                \item \textbf{API/Orchestration Layer:} FastAPI with Pydantic validation, routing, and workflow coordination.
                \item \textbf{Data Layer:} SQLite, YAML registry, and JSON case studies.
            \end{itemize}
        \end{column}
        \begin{column}{0.56\textwidth}
            \centering
            \begin{tikzpicture}[scale=0.78, transform shape,
                node distance=0.55cm,
                layer/.style={rectangle, draw=ntu-blue, fill=ntu-blue!10, text width=5.8cm, minimum height=0.65cm, align=center, font=\small},
                sublayer/.style={rectangle, draw=ntu-blue!50, fill=white, text width=2.6cm, minimum height=0.5cm, align=center, font=\scriptsize},
                arrow/.style={-{Stealth[length=2mm]}, thick, ntu-blue}
            ]
                \node[layer] (pres) {Streamlit Dashboard};
                \node[layer, below=0.4cm of pres] (api) {FastAPI (REST API)};

                \node[sublayer, below left=0.5cm and -0.3cm of api] (eval) {/api/evaluate};
                \node[sublayer, right=0.15cm of eval] (conf) {/api/conflicts};
                \node[sublayer, below right=0.5cm and -0.3cm of api] (par) {/api/pareto};

                \node[layer, below=1.6cm of api] (engine) {Scoring Engine + Harm Assessment};

                \node[layer, below=0.4cm of engine] (data) {YAML Registry | SQLite | Case Study JSON};

                \draw[arrow] (pres) -- node[right, font=\scriptsize]{REST} (api);
                \draw[arrow] (api) -- (eval);
                \draw[arrow] (api) -- (conf);
                \draw[arrow] (api) -- (par);
                \draw[arrow] (engine) -- (data);
                \draw[arrow, dashed] (eval.south) -- ++(0,-0.25) -| (engine.north west);
                \draw[arrow, dashed] (par.south) -- ++(0,-0.25) -| (engine.north east);
            \end{tikzpicture}
        \end{column}
    \end{columns}
\end{frame}
% NOTES: Emphasize clean separation of concerns. The Streamlit frontend communicates
% exclusively through the REST API. The routers (evaluate.py, conflicts.py, pareto.py)
% belong to the API/orchestration layer, handling routing, request validation, and
% workflow coordination. The scoring engine and harm assessment module form the engine
% layer, containing the analysis, scoring, and optimization logic. Note the acknowledged
% technical debt: conflict analysis logic currently resides in the API router layer,
% not in a standalone engine module. Approximately 60 seconds.
% Transition: "The scoring engine at the core uses TOPSIS with a product pooling mechanism."

% ============================================================
% S11: ALGORITHM - TOPSIS SCORING ENGINE
% ============================================================
\begin{frame}{Algorithm: TOPSIS Scoring Engine}
    \textbf{Product Pooling} combines stakeholder and framework weights~\cite{hwang1981multiple, behzadian2012topsis}:
    \[
        w_{\text{eff}} \propto w_s \odot w_f
    \]

    \textbf{TOPSIS Closeness Coefficient} (distance to ideal vs.\ anti-ideal):
    \[
        C = \frac{d^-}{d^+ + d^-} \quad \in [0, 1]
    \]

    \vspace{6pt}
    \begin{itemize}
        \item Single-alternative evaluation uses synthetic best/worst anchors to keep $C$ well-defined.
        \item Risk classification: $C \geq 0.80$ Low, $0.60$--$0.80$ Medium, $0.40$--$0.60$ High, $< 0.40$ Critical.
    \end{itemize}
\end{frame}
% NOTES: Product pooling ensures a dimension scores high only if BOTH the stakeholder
% and the framework consider it important. The synthetic anchor trick is a critical
% adaptation for single-system evaluation. Approximately 90 seconds.
% Anticipated question: "Why TOPSIS over AHP or VIKOR?"
% Response: "TOPSIS is computationally efficient, produces a clear scalar score, and
% performs favorably against VIKOR in comparative analyses. AHP requires pairwise
% comparisons that would burden the user interface."
% Transition: "The conflict detection module uses Spearman rank correlation."

% ============================================================
% S12: ALGORITHM - CONFLICT DETECTION
% ============================================================
\begin{frame}{Algorithm: Conflict Detection (Spearman $\rho$)}
    MEDF computes two types of stakeholder correlation.

    \vspace{6pt}
    \begin{enumerate}
        \item \textbf{Priority Conflict (Weights-Only):} Spearman $\rho$ between raw stakeholder weight vectors. Measures general agreement on dimension importance.
        \item \textbf{System-Salience Conflict (Contribution-Based):} Spearman $\rho$ between contribution vectors (effective weights $\times$ dimension scores). Measures agreement for a specific AI system.
    \end{enumerate}

    \vspace{6pt}
    \textbf{Critical Insight:} The contribution-based metric produces different conflict levels across case studies, even with identical stakeholder weights. This captures case-specific dynamics.
\end{frame}
% NOTES: The distinction between these two metrics is a differentiating feature.
% Contribution-based conflict is what makes the analysis system-aware.
% Approximately 60 seconds.
% Transition: "When conflict is high, MEDF generates Pareto-optimal consensus solutions."

% ============================================================
% S13: ALGORITHM - PARETO AND HARM
% ============================================================
\begin{frame}{Algorithm: Pareto Consensus and Harm Assessment}
    \textbf{NSGA-II} objective function with salience weighting~\cite{deb2002nsga}:
    \[
        f_k(w_c) = \sum_i s_i \cdot |w_{c,i} - w_{k,i}|
    \]

    \textbf{Harm Score} combines system performance and governance disagreement:
    \[
        h = 0.7 \times (1 - s_{\text{norm}}) + 0.3 \times \bar{d}
    \]

    \vspace{4pt}
    \begin{itemize}
        \item The 0.7/0.3 split reflects a design decision: system performance dominates risk classification, with stakeholder disagreement as an aggravating factor.
    \end{itemize}
\end{frame}
% NOTES: The salience factor s_i ensures the optimizer prioritizes dimensions that
% matter most for the system under evaluation. The harm score is a composite metric.
% Approximately 45 seconds.
% Transition: "Let me now present the three case studies used to validate the platform."
